\section{Presentation}
\label{sec:presentation}

This project's aim was designing and developing a multi-page website about Cybersecurity and Digital Privacy.
The work was divided among several topic groups, with the aim of producing an organized website in which the different
sections shared a consistent structure, navigation system, and visual style.

\subsection{Goals}
\label{sec:goals}

The main goal was to create an informative and accessible website that introduces important topics related to cybersecurity and digital privacy.
The website covers different areas, including common cyber threats, authentication, risks introduced by artificial intelligence, privacy and famous incidents, good security practices, and secure communication.

The project also aimed to:
\begin{itemize}

    \item divide a large website into smaller sections that could be developed
    independently by different groups;
    \item explain cybersecurity concepts through dedicated pages, while also doing the required research and attach the sources;
    \item provide consistent navigation between the pages and their topic overview pages;
    \item apply a shared CSS structure so that the website would have a coherent
    visual identity;
    \item practice collaborative development, version control, and integration
    using Git.

\end{itemize}

\subsection{Team and division of work}
\label{sec:team}

The project was developed by a team of 22 students. The work was organized into
four main topic groups:

\begin{itemize}

    \item \textbf{Common Threats and Malware:} led by Adam Balice.
    \item \textbf{Privacy and Famous Incidents:} led by Carlo Augusto Conti.
    \item \textbf{Good Security Practices and Secure Communication:} led by Raul Moreno.
    \item \textbf{Authentication and Risks Brought by AI:} led by Matteo Mignone.
\end{itemize}

Each topic leader coordinated their group and organized the content and HTML pages, ensuring proper integration with the website's navigation.
Angelica Vetere and Vittorio Lo Mele led the project, Giulio Finocchiaro and Lazar Bogosavljevic led the CSS work, and Vittorio managed Git and repository integration.

\subsection{Timeline}
\label{sec:timeline}

Throughout the three weeks allocated to the project, our team successfully reached several important milestones, making steady progress towards the completion of the website.

\begin{itemize}

    \item \textbf{Week 1:} During the first week, we assigned specific roles and responsibilities to each team member to ensure an efficient distribution of tasks. We began outlining the website's overall structure and conducting research on our assigned cybersecurity topics. At the same time, we prioritized the development of the CSS stylesheets to meet our first project deadline and establish a consistent visual design across the website.
    \item \textbf{Week 2:} During the second week, our main focus shifted towards completing the HTML pages and integrating the previously developed CSS stylesheets. We worked on ensuring that the content was properly structured and that the visual design remained consistent throughout the website.
    \item \textbf{Week 3:} In the final week, we focused on refining the website and implementing the remaining adjustments needed for its completion. This included identifying and fixing issues in various sections, improving existing features, and adding supplementary documentation to the Git repository to better organize and document our work.

\end{itemize}

\subsection{Final result}
\label{sec:result}

The result is a website bringing together the work of the four topic groups in one shared project.
Its pages cover a range of cybersecurity and digital privacy subjects, while the shared styling, topic overview pages, and article navigation help users move between related content.

The development process also provided experience in dividing a large task, working collaboratively, using Git, applying shared CSS, reviewing other members' work, and integrating independently developed pages.
Some pages required corrections and additional coordination before the website could be
assembled consistently.
